\section{Discussion and Limitations}
\label{sec:discussion}

Put together, the results describe a short, sharp shock. On the onset day a typical storm makes
objects fall roughly 60\% faster ($-30$ against a pooled quiet mean of about $-50$~m/day). Most of
the response comes near onset, with smaller negative values still visible on days 4--5. Pooled
over days 0--2, NRLMSIS~2.1 shows little mean error, but that average hides two problems. In
ordinary storms it gets the timing wrong and under-predicts decay at peak ap, and in the largest
storms its totals are wrong in storm-specific ways. That matters more for short-term conjunction
screening than for lifetime estimates. Vulnerability follows the drag equation: drag sensitivity
times an altitude factor, with lower objects hit harder per unit drag. And in our exploratory
look at Starlink, the observed orbit changes are consistent with an approximately complete
offset of the predicted loss, almost all during the storm.

The study has clear limits, and we list eight. (1) One year of data gives 20 storms, of which 12
have clean event-study windows, so storm-level intervals are wide and class comparisons have
little power. (2) TLEs are mean-element fits, not true altitudes, and they blur changes faster
than about a day. This is why the extreme storms needed the calm-to-calm check. (3) Question~3
covers a single operator, with catalogue-scaled drag sensitivity and a weak placebo. (4) For six
storms the pre-storm quiet ap was raised by an earlier storm. A storm-free quiet ap would fix
this, and in its place we report the run without those storms. (5) The causal forest's intervals
treat rows as independent. (6) We only evaluate NRLMSIS, and other models such as JB2008 or DTM
may behave differently. (7) The Hampel filter runs on the NRLMSIS~2.1 residual and removes
proportionally more rows at storm onset. Our sensitivity runs show the headlines survive without
it, but the filter could still hide some real storm-peak model error. (8) Intervals are
conditional on the fitted drag sensitivities and, for Question~3, on the catalogue scaling, and
this uncertainty is not carried through. Natural next steps are held-out storms, more of the solar cycle, a
second constellation and other density models.
